% c04 —— 第 8 节 量子化，p014–p015

\section{量子化}

20 世纪，经典力学把物理学基本原理的宝座让给了量子力学，但在此过程中辛几何的重要性反而提高了。

在把 Planck 常数 $\hbar$ 取为 $0$ 的极限下，量子力学（应当）再现经典力学；反过来，对给定的经典系统新引入参数 $\hbar$，使 $\hbar\to0$ 的极限再现原来的系统，这样的操作称为量子化。

经典系统的量子化是否总是存在，并不清楚；即使存在，也未必唯一。

分析力学不仅给出描述经典力学的精炼方法，也为进行量子化提供重要线索。

把用 Lagrange 形式给出的经典系统量子化这一问题，在 Feynman 的学位论文中被提出，并被表述为路径积分理论。

Heisenberg 的矩阵力学把位置、动量等物理量替换为矩阵，把 Poisson 括号替换为算子的交换子，以此进行量子化：
\begin{equation}
    p,q\mapsto\hat p,\hat q,\qquad \{q,p\}=1\mapsto[\hat q,\hat p]=\sqrt{-1}\hbar. \tag{8.1}
\end{equation}

另一方面，Schrödinger 的波动力学与矩阵力学等价，但它与 Hamilton–Jacobi 形式的分析力学密切相关。在 Hamilton–Jacobi 形式中，通过作依赖于时间的正则变换，使 Hamilton 函数为零。刻画这种正则变换的生成函数的方程就是 Hamilton–Jacobi 方程，它是把位置与时间作为独立变量的一阶偏微分方程。它把常微分方程的问题变成偏微分方程的问题，乍看之下给人以把问题变难的印象；但它不仅对通过寻找守恒量来简化问题十分有效，还通过强调粒子运动与波传播的类似性，开辟了通向波动力学之路。

从辛流形的立场重新审视这些理论，就得到几何量子化（geometric quantization）这一概念。特别重要的例子，可举出 Lie 群余伴随轨道的几何量子化，以及 Riemann 面上向量丛的模空间的几何量子化。前者与 Lie 群的表示论密切相关，通向 Kostant 关于 Toda 格子及其量子化的研究。后者与共形场论及 Chern–Simons 理论相联系，在低维拓扑的应用等方面取得了显著成果。

与几何量子化的性质稍有不同的一种量子化，是形变量子化（deformation quantization）。把 $A=C^\infty(M)$ 的结合环的形变（即 $A[[\hbar]]$ 上 $\R[[\hbar]]$ 的结合代数结构 $\star$，使它在模 $\hbar$ 后成为 $A$ 的通常乘法）写作
\begin{equation}
    f\star g=fg+\hbar B_1(f,g)+\hbar^2B_2(f,g)+\cdots \tag{8.2}
\end{equation}
这里假定 $B_i$ 关于两个变元分别是有限阶微分算子。在积 $\star$ 的集合上，作为 $A[[\hbar]]$ 的 $\R[[\hbar]]$ 模的自同构
\begin{equation}
    D\colon f\mapsto f+\hbar D_1(f)+\hbar^2D_2(f)+\cdots \tag{8.3}
\end{equation}
作用如下：
\begin{equation}
    \star\mapsto\star',\qquad a\star' b:=D\left(D^{-1}(a)\star D^{-1}(b)\right) \tag{8.4}
\end{equation}
这里 $D_i$，$i=1,2,\dots$ 是有限阶微分算子。

\begin{exercise}
证明：$\star$ 的结合律
\begin{equation}
    (f\star g)\star h=f\star(g\star h) \tag{8.5}
\end{equation}
在 $\hbar$ 的一次项给出
\begin{equation}
    fB_1(g,h)-B_1(fg,h)+B_1(f,gh)-B_1(f,g)h=0 \tag{8.6}
\end{equation}
\end{exercise}

\begin{exercise}
证明：自同构 (8.3) 作用在 $B_1$ 上为
\begin{equation}
    B_1'(f,g)=B_1(f,g)-fD_1(g)+D_1(fg)-D_1(f)g \tag{8.7}
\end{equation}
\end{exercise}

\begin{exercise}
证明：把 $B_1$ 分成对称部分与反对称部分
\begin{align}
    B_1^+(f,g)&=\frac{1}{2}\left(B_1(f,g)+B_1(g,f)\right), \tag{8.8}\\
    B_1^-(f,g)&=\frac{1}{2}\left(B_1(f,g)-B_1(g,f)\right), \tag{8.9}\\
    B_1^+(f,g)&=B_1^+(f,g)+B_1^1(g,f) \tag{8.10}
\end{align}
时，自同构 (8.3) 作用为
\begin{align}
    B_1'^{\,+}(f,g)&=B_1^+(f,g)-fD_1(g)+D_1(fg)-D_1(f)g, \tag{8.11}\\
    B_1'^{\,-}(f,g)&=B_1^-(f,g) \tag{8.12}
\end{align}
\end{exercise}

% [Tutor Note: 原文 (8.10) 右端第二项的上标印作 1，即 B_1^1(g,f)，与 (8.8)(8.9) 定义 B_1^+、B_1^- 的记号不相符，疑为笔误；此处照原文保留，请审校确认。]

对给定的 $\star$ 积，定义 $A$ 的 Poisson 括号为
\begin{equation}
    \{f,g\}:=\lim_{\hbar\to0}\frac{1}{\hbar}[f,g]=2B_1^-(f,g) \tag{8.13}
\end{equation}
其中
\begin{equation}
    [f,g]:=f\star g-g\star f \tag{8.14}
\end{equation}
是关于 $\star$ 积的交换子。

\begin{definition}
交换环 $A$ 与映射 $\{-,-\}\colon A\times A\to A$ 的组 $(A,\{-,-\})$ 称为 Poisson 代数（Poisson algebra），是指 $A$ 满足
\begin{enumerate}
    \item 反对称性：$\{a,b\}=-\{b,a\}$
    \item Jacobi 恒等式：$\{\{a,b\},c\}+\{\{a,b\},c\}+\{\{a,b\},c\}=0$
    \item Leibniz 法则：$\{a,bc\}=\{a,b\}c+b\{a,c\}$
\end{enumerate}
\end{definition}

% [Tutor Note: 原文 定义8.1 的 Jacobi 恒等式三项被印成完全相同的 {{a,b},c}，按通常写法应为 {{a,b},c}+{{b,c},a}+{{c,a},b}=0；此处照原文保留，请审校确认。]

\begin{exercise}
证明：由 $A$ 与 $A[[\hbar]]$ 的 $\star$ 积按 (8.13) 定出的 Poisson 括号 $\{-,-\}$，其组 $(A,\{-,-\})$ 构成 Poisson 代数。（提示：看由 $\star$ 的结合律导出的交换子积 $[-, -]$ 的 Jacobi 恒等式与 Leibniz 法则中关于 $\hbar$ 的首项。）
\end{exercise}

% CHUNK c04 DONE pages p014–p015
